\begin{resumen}
La transformación digital del sector de la salud exige el desarrollo de infraestructuras de datos soberanas 
que respondan a las particularidades epidemiológicas y demográficas de cada región. En el ámbito de la 
dermatología computacional, los principales repositorios internacionales de imágenes ---como ISIC, 
HAM10000, BCN20000, PAD-UFES-20 y Derm7pt--- exhiben sesgos demográficos sistemáticos, con una 
desproporción a favor de fototipos claros (Fitzpatrick I--III), lo que compromete 
la equidad algorítmica y la fiabilidad de los modelos de inteligencia artificial cuando se aplican a 
poblaciones con alta diversidad fenotípica, como la cubana.

La presente investigación aborda el diseño, desarrollo e implementación de la Base de Imágenes 
Dermatoscópicas del ecosistema DermaUH, una plataforma integral concebida como el núcleo de persistencia 
y análisis del ecosistema informático liderado por la Facultad de Matemática y Computación de la 
Universidad de La Habana. El sistema propuesto adopta una arquitectura por capas fundamentada en 
los principios de la Arquitectura Limpia, materializada sobre la plataforma .NET~10, el lenguaje 
C\#~13, el motor relacional PostgreSQL~17 y la interfaz interactiva Blazor Server. El esquema de 
datos diseñado incorpora más de 30 atributos clínicos estructurados por registro, incluyendo 
variables demográficas, antecedentes epidemiológicos, indicadores histológicos y metadatos 
institucionales.
Entre las contribuciones principales se destacan: la implementación de un motor de validación clínica con 
reglas condicionales que garantiza la coherencia médica de los metadatos; un sistema de control de acceso 
basado en roles con autenticación JWT y OAuth~2.0; mecanismos de exportación masiva en formatos CSV y ZIP 
orientados a la generación de conjuntos de datos listos para el aprendizaje profundo; y un cuadro de mando 
analítico con visualizaciones estadísticas demográficas y geográficas. La plataforma resultante constituye 
un activo tecnológico soberano, diseñado para capturar la diversidad fenotípica de la población cubana y 
contribuir a la reducción de la brecha de equidad algorítmica documentada en la literatura científica 
internacional.

\textbf{Palabras clave:} base de imágenes dermatoscópicas, ecosistema DermaUH, soberanía tecnológica, 
inteligencia artificial, equidad algorítmica, arquitectura de software, metadatos clínicos.
\end{resumen}

\begin{abstract}
The digital transformation of the healthcare sector demands the development of sovereign data 
infrastructures that address the epidemiological and demographic particularities of each region. 
In the field of computational dermatology, leading international image repositories ---such as ISIC, 
HAM10000, BCN20000, PAD-UFES-20, and Derm7pt--- exhibit systematic demographic biases, with a 
disproportion in favor of light skin phototypes (Fitzpatrick I--III), thereby compromising algorithmic 
fairness and the reliability of artificial intelligence models when applied to populations with high 
phenotypic diversity, such as the Cuban population.

This research addresses the design, development, and implementation of the Dermoscopic Image Database for 
the DermaUH ecosystem, a comprehensive platform conceived as the core persistence and analysis component 
of the informatics ecosystem led by the Faculty of Mathematics and Computer Science at the University of 
Havana. The proposed system adopts a layered architecture grounded in Clean Architecture principles, built 
upon the .NET~10 platform, the C\#~13 programming language, the PostgreSQL~17 relational engine, and the 
Blazor Server interactive interface. The data schema incorporates over 30 structured clinical attributes 
per record,including demographic variables, epidemiological history, histological indicators, and 
institutional metadata.

Key contributions include: the implementation of a clinical validation engine with conditional rules that 
ensures the medical consistency of metadata; a role-based access control system with JWT and OAuth~2.0 
authentication; bulk export mechanisms in CSV and ZIP formats oriented toward generating AI-ready datasets; 
and an analytical dashboard featuring demographic and geographic statistical visualizations. The resulting 
platform constitutes a sovereign technological asset, designed to capture the phenotypic diversity of the 
Cuban population and contribute to reducing the algorithmic fairness gap documented in the international 
scientific literature.

\textbf{Keywords:} dermoscopic image database, DermaUH ecosystem, technological sovereignty, artificial 
intelligence, algorithmic fairness, software architecture, clinical metadata.
\end{abstract}